\begin{abstract}
Satellite altimetry is widely used to estimate global mean sea level change (GMSLC), but the standard approach treats sea surface height change (SSHC) as a proxy for sea level change (SLC), neglecting the gravitationally self-consistent redistribution of water that links ice mass loss to the sea surface. This report quantifies the systematic bias introduced by this proxy and develops a principled Bayesian alternative.

Error quantification for the three major ice sheets shows that the SSHC proxy introduces source-dependent biases: GMSLC is underestimated by approximately 5.8--8.2\% for Greenland melt and 0.5--2\% for East Antarctic melt, and overestimated by approximately 1--2\% for West Antarctic melt. Because the bias depends on where ice is melting, no scalar correction can account for it; correcting for the error requires knowledge of the very quantity being estimated.

In place of the proxy, we apply the infinite-dimensional Bayesian inversion framework of \textcite{heathcoteScaleableBayesianFramework2026}, implemented in \citetitle{al-attarPygeoinfPythonLibrary2025}, to recover ice thickness change directly from SSH altimetry. Priors are defined as Gaussian measures over Sobolev spaces on the sphere, encoding physically motivated beliefs about the spatial structure of mass change via a novel Richards curve construction conditioned on ice thickness. The posterior is obtained in closed form and propagated through the full sea level physics without discretisation artefacts. Across 261 synthetic twin experiments, 60.5\% of true GMSLC values fall within one posterior standard deviation, compared to 14.6\% under the proxy method, providing direct empirical evidence of well-calibrated uncertainty quantification.

The inversion is extended to a joint inference over ice thickness change, firn compaction, and ocean dynamic topography, combining SSH altimetry, ice altimetry, and GRACE gravity observations in a single composite model space. The joint forward operator is factorised to reduce computational cost by a factor of three without approximation. Knock-out tests demonstrate that each data source contributes non-redundant information, and the posterior successfully disentangles contributions that are entangled at the level of the observations. Sensitivity tests confirm robustness to prior mis-specification in the inversion results, though computational efficiency is sensitive to the prior covariance amplitude.

Taken together, the methods developed here offer three technical advances over both the SSHC proxy and existing discrete parameterisation approaches: principled uncertainty propagation through the full forward physics; spatially correlated posteriors whose structure is inherited from physically motivated priors rather than imposed by a discretisation grid; and joint inference over competing signals with formal uncertainty quantification on the trade-offs between them.
\end{abstract}
